\documentclass[a4paper, 12pt]{article}

%%% SST LAB PROTOCOLL PREAMBLE
%%% 2019
%%%%%%%%%%%%%%%%%%%%%%%%%%%%%%%


%%% PACKAGES
%%%%%%%%%%%%%%%%%%%%%%%%%%%

\usepackage[ngerman]{babel}

\usepackage[utf8]{inputenc}
\usepackage{amsmath}
\usepackage{pgfplots}
\usepackage{tikz}
\usepackage[many]{tcolorbox}
\usepackage{graphicx}
\graphicspath{ {./graphics/} }
\usepackage{pdfpages}
\usepackage{dashrule}
\usepackage{float}
\usepackage{siunitx}
\usepackage{trfsigns}
\usepackage{booktabs}
\usepackage[european]{circuitikz}
\usepackage{tcolorbox}

%%% DOCUMENT GEOMETRY
%%%%%%%%%%%%%%%%%%%%%%%%%%%

\usepackage{geometry}
\geometry{
 a4paper,
 total={0.6180339887498948\paperwidth,0.6180339887498948\paperheight},
 top = 0.1458980337503154\paperheight,
 bottom = 0.1458980337503154\paperheight
 }
\setlength{\jot}{0.013155617496424828\paperheight}
\linespread{1.1458980337503154}

\setlength{\parskip}{0.013155617496424828\paperheight} % paragraph spacing


%%% COLORS
%%%%%%%%%%%%%%%%%%%%%%%%%%%

\definecolor{red1}{HTML}{f38181}
\definecolor{yellow1}{HTML}{fce38a}
\definecolor{green1}{HTML}{95e1d3}
\definecolor{blue1}{HTML}{66bfbf}
\definecolor{hsblue}{HTML}{00b1db}
\definecolor{hsgrey}{HTML}{afafaf}

%%% CONSTANTS
%%%%%%%%%%%%%%%%%%%%%%%%%%%
\newlength{\smallvert}
\setlength{\smallvert}{0.0131556\paperheight}


%%% COMMANDS
%%%%%%%%%%%%%%%%%%%%%%%%%%%

% differential d
\newcommand*\dif{\mathop{}\!\mathrm{d}}

% horizontal line
\newcommand{\holine}[1]{
  	\begin{center}
	  	\noindent{\color{hsgrey}\hdashrule[0ex]{#1}{1pt}{3mm}}\\%[0.0131556\paperheight]
  	\end{center}
}

% mini section
\newcommand{\minisec}[1]{ \noindent\underline{\textit {#1} } \\}

% quick function plot
\newcommand{\plotfun}[3]{
  \vspace{0.021286\paperheight}
  \begin{center}
    \begin{tikzpicture}
      \begin{axis}[
        axis x line=center,
        axis y line=center,
        ]
        \addplot[draw=red1][domain=#2:#3]{#1};
      \end{axis}
    \end{tikzpicture}
  \end{center}
}

% box for notes
\newcommand{\notebox}[1]{

\tcbset{colback=white,colframe=green1!100!black,title=Note!,width=0.618\paperwidth,arc=0pt}

 \begin{center}
  \begin{tcolorbox}[]
   #1 
  \end{tcolorbox}
 
 \end{center} 
 
}

% box for equation
\newcommand{\eqbox}[2]{
	
	\tcbset{colback=white,colframe=green1!100!black,title=,width=#2,arc=0pt}
	
	\begin{center}
		\begin{tcolorbox}[ams align*]
				#1
		\end{tcolorbox}
		
	\end{center} 
	
}
% END OF PREAMBLE

%%%%%%%%%%%%%%%%%%%%%%%%%%%%%%%%%%%%%

\begin{document}

%%%%%%%%%%%%%%%%%%%%%%%%%%%%%%%%%%%%%
  \includepdf{./titlepage/titlepage.pdf}
  \clearpage
  \setcounter{page}{1}
%%%%%%%%%%%%%%%%%%%%%%%%%%%%%%%%%%%%%
  \section{Timer A Initialisierung}
  Timer A wird hauptsächlich über das \inlinecode{TACTL} Register initialisert.
  Dort kann u.a. die Taktquelle \inlinecode{TASSEL} ausgewählt werden und es können Timermodus
  \inlinecode{MC}, Frequenzteiler \inlinecode{ID} und Interrupt Enable \inlinecode{IE} gesetzt werden.
  Entsprechend des Timermodus muss dann möglicherweise noch das \inlinecode{TACCR0} Register
  gesetzt werden.
  
  \section{Timer A Start}
  Der Timer wird über die \emph{Mode Control} (\inlinecode{MC}) Bits in einen Modus, d.h.
  up-, up-down-, bzw. continuous-Mode versetzt und dadurch gestartet. Durch
  löschen dieser Bits wird der Timer gestoppt. 

  \section{Output Mode und Zählmodus}
  Zur Frequenzerzeugung mit dem Timer sollte dieser im \emph{up-Mode} Zählmodus
  und \emph{Reset/Set} Output Mode sein. Der Wert in \inlinecode{TACCR0} ist die
  Anzahl der Zählungen des Timers bis zum Reset auf 0 und bestimmt damit die Periodendauer $T$
  \[ \inlinecode{TACCR0} = T \cdot f_{CLK} - 1\]

  Analog wird der Wert der Impulsbreite $\tau$ durch den Wert in
  \inlinecode{TACCRn} bestimmt 
  \[\inlinecode{TACCRn} = \tau \cdot f_{CLK} - 1\]

  Dabei ist es wichtig, auf die richtigen Einheiten zu achten. Unter Umständen
  werden noch Konvertierungsfaktoren benötigt.

  \section{Minimale und maximale Frequenz}
  Die minimale Frequenz die erzeugt werden kann, ist
  \[f_{min} = \frac{f_{CLK}/d}{2\cdot 65536}\]
  wenn \inlinecode{TACCR0} auf den
  Maximalwert (\inlinecode{0xFFFFFF}) gesetzt, der Taktteiler auf dem Maximalwert
  ($d=8$) und der Timer im up-down-Mode ist.\\

  Die maximale Frequenz, die erzeugt werden kann, ist die halbe
  Timertaktfrequenz im up-Mode
  \[f_{max} = \frac{f_{CLK}}{2}\]
  


 \end{document}

% Local Variables:
% TeX-engine: luatex
% End:
